\section{Experiments}

Our principal comparison holds the backbone and adaptation status fixed: all systems in Table~\ref{tab:frozen-results} use frozen base models without benchmark-specific training. We separately use a shared-checkpoint adaptation study to diagnose proposal-interface effects; it is not mixed into the frozen comparison.

\paragraph{Data and Protocol.}
\bench{} contains 3,500 L3 development scenarios and 500 template-disjoint L4 test scenarios, covering 15,203 memory records, 28,409 chronological events, and seven failure modes. Prompts, corruption operators, and calibration constants are fixed on L3. L4 is used once for final reporting. Every system must produce the same five-field response contract and is scored by the same deterministic evaluator. Invalid parses and schema violations are retained as failures rather than silently repaired.

\paragraph{Systems.}
The five frozen baselines are: \emph{Frozen Direct Prompting}, which predicts the final response from the task view; \emph{Full Context}, which exposes the complete public scenario; \emph{Lexical RAG}, which retrieves evidence by lexical relevance; \emph{Time-Aware RAG}, which adds recency-aware ranking; and \emph{Summary Memory}, which compresses prior events. \method{} uses the same frozen backbone to propose typed patch actions, followed by the Revision Guard and canonical projection. Thus the comparison changes the inference interface and mediator, not model weights.

\begin{table*}[t]
\centering
\small
\setlength{\tabcolsep}{5pt}
\begin{tabular}{lrrrrrrrr}
\toprule
& \multicolumn{4}{c}{Memory-State Accuracy (\%)} & \multicolumn{4}{c}{Schema Compliance (\%)} \\
\cmidrule(lr){2-5}\cmidrule(lr){6-9}
System & Qwen & Phi-4 & Mistral & Avg. & Qwen & Phi-4 & Mistral & Avg. \\
\midrule
Frozen Direct Prompting & 82.2 & 74.0 & 0.0 & 52.1 & 100.0 & 59.6 & 0.0 & 53.2 \\
Full Context & 74.6 & 74.1 & 2.9 & 50.5 & 100.0 & 63.8 & 4.2 & 56.0 \\
Lexical RAG & 79.9 & 79.2 & 9.3 & 56.1 & 90.6 & 58.2 & 1.8 & 50.2 \\
Time-Aware RAG & 76.4 & \textbf{85.1} & 6.5 & 56.0 & 98.4 & 63.0 & 11.0 & 57.5 \\
Summary Memory & 81.1 & 81.3 & 0.0 & 54.1 & 76.2 & 54.4 & 0.0 & 43.5 \\
\method{} & \textbf{83.9} & 82.9 & \textbf{84.7} & \textbf{83.8} & \textbf{100.0} & \textbf{100.0} & \textbf{100.0} & \textbf{100.0} \\
\bottomrule
\end{tabular}
\caption{Matched frozen-backbone comparison on the 500-case L4 test set. No row uses LoRA or benchmark-specific adaptation. We emphasize memory-state accuracy, the benchmark's persistent-state target, and schema compliance, which determines whether an output can be consumed without repair. Full diagnostic vectors are reported separately.}
\label{tab:frozen-results}
\end{table*}

\paragraph{Persistent-State Reliability.}
Across the three completed frozen backbones, \method{} raises average memory-state accuracy from the strongest baseline average of $56.1\%$ to $83.8\%$. The largest difference occurs on Mistral-Nemo: the strongest baseline reaches only $9.3\%$, while typed patches followed by deterministic projection reach $84.7\%$. This is not merely an average dominated by one model: \method{} is best on Qwen and Mistral and remains competitive with the strongest Phi-4 baseline ($82.9\%$ versus $85.1\%$).

\paragraph{Interface Robustness.}
All \method{} runs satisfy the output schema, whereas baseline compliance ranges from $0\%$ to $100\%$ and collapses on Mistral. The result isolates a practical role for mediation: model-specific surface forms are normalized before persistent state is committed. The Guard cannot make an omitted proposal correct, but it can prevent malformed output from becoming an implicit write protocol.

\paragraph{Diagnostic Trade-offs.}
The full vectors prevent overclaiming. \method{} does not dominate every response channel. On Qwen, for example, frozen direct prompting has higher decision $F_1$ and evidence $F_1$, while \method{} has higher memory-state accuracy and complete schema compliance. On Phi-4, \method{} obtains the strongest diagnosis accuracy ($37.4\%$) but trails Time-Aware RAG in decision $F_1$. These differences support the benchmark design: answer-level, evidence-level, and persistent-state correctness are related but not interchangeable.

\paragraph{Why Joint Is Secondary.}
Joint revision success requires every scored component to be correct for the same case. It is useful as a strict complete-case statistic, but a zero does not distinguish malformed output, a wrong diagnosis, or a single evidence mismatch from a wrong persistent state. We therefore retain Joint in reproducibility artifacts rather than using it as the headline measure.

\subsection{Mediator and Adaptation Diagnostics}

\paragraph{Guard Effect.}
For frozen Qwen, Phi-4, and Mistral, the Guarded path improves memory-state accuracy over direct projection of the same typed-action proposal by $4.9$, $3.2$, and $0.8$ percentage points, respectively. The small or zero differences on other channels are expected: the Guard authorizes and projects proposed transitions; it is not a second semantic predictor. Controlled invalid-action injection is the appropriate setting for measuring rejection and unauthorized write-through.

\paragraph{Shared-Checkpoint Adaptation.}
An auxiliary QLoRA experiment trains one multitask adapter per backbone on both final-state and patch-action targets. Checkpoints are saved every 128 steps through step 1,024 and selected solely by minimum L3 validation loss. Final-State Control and \method{} then use the exact same selected adapter, preventing interface-specific checkpoint selection. On the completed Qwen run, Final-State Control achieves higher complete-case accuracy, while \method{} yields explicit action validity and per-transition audit coverage. We treat this as an interface trade-off, not as evidence that adaptation is required by the runtime.

\paragraph{Natural versus Injected Safety.}
The natural L4 runs contain no measured stale-write-through events under the current stale-reuse statistic, so they cannot establish a reduction in that rate. Safety claims instead come from the Guard's conditional construction and must be complemented by controlled corruption tests reporting invalid-action rejection, replay consistency, and unauthorized write-through.
